\section{Conclusion}
\label{sec:conclusion}
Repairing garbled text in generated videos is not the same problem as editing text: the source is a near miss, it is often already right, and everything around it should survive. \method treats it as a sequence of evidence-gated decisions---leave the clip alone, lay the text out on the old glyphs or on the tracked box, transport the old appearance or not---around a rendering step whose spelling cannot drift. This makes it more reliable than the strongest generative editor on our benchmark and confines its changes to the text, but its glyphs still look rendered rather than generated. Our analysis and two pre-registered attempts indicate what closing that gap without giving up correctness requires: letter segmentation and matching reliable enough to reuse the characters that the generator already drew correctly, an appearance measure that sees glyph shape and orientation, and generative edits that are accepted only when they can be verified.
